% Chapter 7 -- Rubric: PO12 (12.1)
\chapter{Conclusion}
\label{ch:conclusion}

\section{Conclusion and Future Work}
\label{sec:conclusion}

\subsection{Achievement against objectives}

\begin{table}[htbp]
  \centering
  \caption{Achievement against the objectives of \Cref{sec:objectives}}
  \label{tab:objectives-status}
  \begin{tabularx}{\textwidth}{>{\hsize=0.22\hsize}Y>{\hsize=0.45\hsize}Y>{\hsize=2.33\hsize}Y}
    \toprule
    Obj. & Status & Basis \\
    \midrule
    O1 & Achieved & 1{,}688 skills, 1{,}743 edges, all 183 syllabus topics populated; validated acyclic, depth 9, zero isolated skills, zero unresolved references \\
    O2 & Achieved & Bloom-conditioned BKT implemented from the equations; exercised on every response by the verification harness \\
    O3 & Achieved in implementation & Ancestor pull-up and conjunctive gating implemented; the ancestor-ordering invariant is checked explicitly and holds. Whether the resulting estimates match real learner knowledge is untested \\
    O4 & Achieved & A 30-question diagnostic completes and moves mastery on 43 skills --- more skills than questions asked, which is the propagation doing its work \\
    O5 & Achieved & Topic practice runs to the mastery threshold; spillover fires on the intended evidence \\
    O6 & \textbf{Partially achieved} & Pipeline built, validated, and producing verified items; coverage of the full specification space incomplete \\
    O7 & Achieved & Seven-route web client with diagnostic, practice and mastery views, bilingual, with Bangla and \LaTeX{} rendering \\
    O8 & Achieved & 18-check end-to-end harness, credential-free, passing completely \\
    \bottomrule
  \end{tabularx}
\end{table}

The project set out to build a tutor that could tell a student \emph{which}
foundation is missing rather than merely that they are wrong, and the mechanism
for that --- an interpretable per-skill model, propagating along an explicit
prerequisite graph, gated conjunctively --- is built and demonstrably
self-consistent. The single most valuable thing the project learned is not in
the engine at all: it is that the knowledge structure, not the inference
algorithm, is the binding constraint on a system like this. A sophisticated
update rule over a graph with 1.1 edges per skill has almost nothing to reason
with. Most of the engineering effort in the second half of the project went into
the ontology for exactly that reason, and that is the finding most likely to
transfer to anyone attempting similar work.

\subsection{Limitations}

Stated plainly, and in the order that matters:

\begin{enumerate}
  \item \textbf{No efficacy evidence.} The system has never been used by a real
    learner. Every claim about learning benefit is an argument from design.
  \item \textbf{No fitted parameters.} All model parameters are pedagogical
    priors; the literature is clear that fitted and individualised parameters do
    better \cite{yudelson2013individualized,baker2008more}.
  \item \textbf{The validated ontology is not the served ontology.} The
    application still runs on the 430-skill legacy graph, because adopting the
    rebuild without regenerating content would leave most skills unquestionable.
  \item \textbf{Question-bank coverage is partial.}
  \item \textbf{The ontology has not been reviewed by domain experts} --- 32
    disconnected fragments and a body of hand-authored chemistry edges await a
    chemist and a physicist.
  \item \textbf{Production-readiness gaps:} no password authentication,
    in-memory session state, wide-open CORS, no accessibility audit, no
    performance benchmark.
\end{enumerate}

\subsection{Future work}
\label{sec:future}

\textbf{Immediate, and blocking deployment.} Password authentication and
server-side sessions; persistent session state; CORS restriction. These are
small, well-understood, and are the difference between a demonstrator and
something that may lawfully hold real learner data.

\textbf{Completing the knowledge base.} A budgeted generation run against the
$1{,}688 \times 6$ specification space, followed by adoption of the rebuilt
ontology in the correct order (generate, then switch); domain-expert review of
the 32 fragments and the authored edges.

\textbf{Establishing efficacy --- the most important item.} A controlled study
with real candidates: does the system's mastery estimate predict performance on
held-out questions better than a baseline, measured by AUC as the knowledge
tracing literature does; and does directing study by the system's diagnosis
produce better outcomes than self-directed study? Until this is done, the central
claim of the project is unverified.

\textbf{Model improvements once data exists.} Fitted and individualised BKT
parameters \cite{yudelson2013individualized}; contextual estimation of slip and
guess \cite{baker2008more}; and, with sufficient data, a comparison against a
deep knowledge tracing baseline \cite{piech2015deep}. The engine was
deliberately built with a stateless, replaceable model interface so that this
substitution requires no change to its call sites.

\textbf{Pedagogical extensions.} Forgetting/decay in the mastery estimate, which
the current model omits entirely and which matters over a months-long
preparation period \cite{ebbinghaus1885memory}; spaced review scheduling; and
session pacing that resists the intensification risk noted in
\Cref{sec:society}.

\section{Continuous Engagement and Self-directed Learning}
\label{sec:lessons}

Little of what this project required was covered by prior coursework, and the
learning it forced was mostly self-directed.

\textbf{Domain literature learned independently.} The knowledge-tracing
literature was new to the team: the original BKT formulation
\cite{corbett1995knowledge}, the variants that refine its noise parameters and
individualise them \cite{baker2008more,yudelson2013individualized,pardos2011kt},
the competing model families \cite{cen2006learning,pavlik2009performance,piech2015deep},
the prerequisite-structured student models closest to this design
\cite{kaser2017dynamic}, and the older knowledge-space theory that supplies the
graph formulation \cite{doignon1985spaces,falmagne2011learning}. The educational
side --- the revised Bloom taxonomy \cite{anderson2001revised} and the zone of
proximal development \cite{vygotsky1978mind} --- was read as source material for
design decisions rather than as background, since both appear directly in the
implemented policy.

\textbf{Emerging technical practice.} Retrieval-augmented generation
\cite{lewis2020retrieval} and multilingual sentence embeddings
\cite{reimers2019sentence} were adopted from current practice; the survey
literature on automatic question generation \cite{kurdi2020systematic} was what
convinced the team that generation without an independent verification stage
would not produce usable assessment items, which directly shaped
\Cref{sec:generation-pipeline}.

\textbf{Tools and technologies learned during the project.} FastAPI and Pydantic;
React with Vite and Tailwind; Supabase and its constraint model; KaTeX and the
specific difficulties of mixing Bangla text with \LaTeX{} mathematics; local and
hosted large-language-model inference including quantised local serving, and the
practical engineering around hosted APIs --- rate limiting, key pooling, adaptive
backoff and resumable batch jobs.

\textbf{Learning from failure, which was the majority of it.} Three episodes
taught more than any tutorial. Measuring local inference before committing to it
would have saved weeks --- the lesson being to measure the throughput of a
critical dependency before building on it. The silent JSON escape corruption
taught that a parser accepting input is not evidence the input was understood,
and that data corruption which raises no error is the expensive kind. The two
over-flagging validators taught that a heuristic must be calibrated against hand
review before its output is treated as a defect list.

\textbf{Practices adopted.} The habit that most changed how the team works is
documentary: writing down the reasoning behind a decision at the moment it is
taken, not afterwards. On a five-person project spanning months, the recurring
bottleneck was not writing code but recovering context --- why a constant has the
value it has, why an approach was abandoned, what a validator's output actually
means. The chain of handover documents in the repository exists because of that,
and it is the practice the team expects to carry into work beyond this course.
